% ch03-foundations --- temporal point processes: the machinery every model chapter instantiates

\section{One machinery, many models}\label{ch03:sec:overview}

Every probabilistic model in this document is an instance of one object: a
counting process characterized by its conditional intensity. The intensity
induces the likelihood we fit, its compensator induces the residuals we
check, and the same calculus justifies the simulators. This chapter builds
that machinery once, so that
\cref{ch:poisson,ch:hawkes,ch:censoring,ch:gof,ch:ranking} can instantiate
\cref{thm:ch03:jacod} with their own $\lam$ and move on.

Five results carry the weight. First, the point-process likelihood
(\cref{thm:ch03:jacod}): one formula, an event term plus an exposure term,
that every fitter in \modrepo{} implements. Second, the ground/mark
factorization (\cref{thm:ch03:groundmark}): the exact identity, promised in
\cref{ch01:sec:factorization}, that splits the marked deal process into a
timing factor and an identity factor and thereby licenses the repository's
two-layer sector$\to$firm architecture. Third, martingale estimation theory
(\cref{ch03:sec:estimation}): why maximum likelihood works for dependent
event data, and which standard errors to trust when the model is wrong.
Fourth, the time-rescaling theorem (\cref{thm:ch03:rescaling}): the
goodness-of-fit engine of \cref{ch:gof}. Fifth, simulation correctness
(\cref{prop:ch03:thinning}): why the generators planned for the synthetic
world of \cref{ch02:sec:synthetic} produce the processes they claim to.
The chapter closes, as \cref{ch:data}'s decision D2.1 requires, with an
explicit account of which data defects (C1)--(C6) this machinery absorbs
and which it cannot.

\section{Counting processes, histories, and conditional intensity}\label{ch03:sec:counting}

\begin{definition}[Counting process]\label{def:ch03:counting}
On the filtered probability space $(\Omega,\F,\{\Ft\},\Prob)$ of
\cref{def:ch02:market}, a \emph{counting process} $N$ is an adapted,
right-continuous process with $N(0)=0$, nondecreasing piecewise-constant
paths, jumps of size $+1$, and $N(t)<\infty$ a.s.\ for every $t \le \Tend$.
A \emph{multivariate counting process} is a finite family
$(N_e)_{e \in \mathcal{E}}$ of counting processes, no two of which jump at
the same time.
\end{definition}

The index $e$ will range over sectors, firms, or firm--transition pairs as
each chapter requires; the filtration $\Ft$ contains the internal history
$\Hist_t$ of the marked deal process \emph{and} the observed covariate
paths, as fixed in \cref{def:ch02:market}. Two pieces of structure make the
theory work. We name both precisely, because a violation of either shows up
first as a software bug.

\paragraph{Predictability.}
A process $H$ is \emph{predictable} if it is measurable with respect to the
$\sigma$-algebra on $\Omega \times [0,\Tend]$ generated by left-continuous
adapted processes. The operational content: the value $H(t)$ must be
computable from information available strictly before $t$, i.e.\ from
$\Ftm$. Left-continuity in $t$ is the working test --- an adapted process
with left-continuous paths is predictable, and every intensity model in
this document is built from such pieces (LOCF covariates held constant
between deals, kernel sums evaluated as left limits, at-risk indicators of
the form $\ind{\entry_i \le t < U_i}$ which are left-continuous in $t$).

\begin{warning}[Predictability is point-in-time discipline]\label{warn:ch03:lookahead}
Predictability of $\lam$ is the continuous-time statement of
\cref{prot:ch02:pit}. If any input to $\lam(t)$ carries a timestamp not
strictly before $t$ --- a covariate as-of joined on reference date instead
of publish date, a kernel sum that includes the event at $t$ itself, a
risk-set flag closed using information revealed at $t$ --- then $\lam$ is
not predictable, $N - \int \atrisk\lam$ is not a martingale with respect to
the filtration actually used, and everything downstream fails
\emph{silently}: the score of \cref{prop:ch03:score} acquires a bias, the
information identities break, and held-out likelihoods flatter the model.
This is why \modrepo{eventtime/likelihood.py} decays its excitation state
to \emph{just before} each event when evaluating $\log\lam$ at that event.
\end{warning}

\paragraph{Compensators.}
The second piece of structure is the decomposition of a counting process
into signal and noise.

\begin{theorem}[Doob--Meyer decomposition for counting processes; stated, not proved]\label{thm:ch03:doobmeyer}
Let $N$ be a counting process with $\E N(\Tend) < \infty$. There exists an
a.s.\ unique predictable, nondecreasing process $\Lam$ with $\Lam(0)=0$,
the \emph{compensator}, such that $M = N - \Lam$ is a martingale. Under
local integrability the same holds with ``local martingale''
\citep[Ch.~II]{abgk1993}.
\end{theorem}

\Cref{def:ch02:intensity} specializes this to the absolutely continuous
case that this document lives in: a \emph{conditional intensity} for $N_e$
is a nonnegative predictable process $\lam_e$ such that
\begin{equation}\label{eq:ch03:compensator}
M_e(t) \;=\; N_e(t) - \int_0^{t} \atrisk_e(u)\,\lam_e(u)\dd u
\end{equation}
is a local martingale, with $\atrisk_e$ the at-risk indicator. The
heuristic that gives the definition its meaning, and that every proof in
this chapter leans on, is the infinitesimal-bin reading of the martingale
property:
\begin{equation}\label{eq:ch03:heuristic}
\lam_e(t)\dd t
 \;=\;
 \Prob\big(N_e \text{ jumps in } [t, t+\dd t) \;\big|\; \Ftm\big)
 \qquad \text{on } \{\atrisk_e(t) = 1\}.
\end{equation}
The intensity is the instantaneous event rate given \emph{everything known
just before now}: internal history, covariate paths, risk-set composition.
It is a modeling object (we write $\lam$ down) and simultaneously a
probabilistic one (\cref{eq:ch03:compensator} pins down what writing it
down commits us to).

\begin{assumption}[Standing regularity]\label{asm:ch03:regularity}
Throughout this chapter $(N_e)_{e\in\mathcal E}$ is a multivariate counting
process such that (i) each $N_e$ admits the intensity
$\atrisk_e\lam_e(\cdot;\thetamodel)$ of \cref{eq:ch03:compensator}, with
$\lam_e$ predictable, locally bounded, and strictly positive at the jump
times of $N_e$; (ii) $\E \int_0^{\Tend}\atrisk_e\lam_e\dd u < \infty$;
(iii) $\lam_e(t;\cdot)$ is twice continuously differentiable in
$\thetamodel$ on a neighborhood of the truth, with predictable, locally
bounded derivative processes, and differentiation under the integral sign
is valid; (iv) no two components jump simultaneously. All models fitted in
this repository satisfy (i)--(iv) by construction (exponential-family
intensities, piecewise-constant covariates, finite windows).
\end{assumption}

\section{The point-process likelihood}\label{ch03:sec:likelihood}

\begin{theorem}[Point-process (Jacod) likelihood]\label{thm:ch03:jacod}
Under Assumption~\ref{asm:ch03:regularity}, let the data on $[0,\Tend]$ be the event
times and types $(t_m, e_m)$, $m = 1,\dots,n$. Then, up to an additive term
not depending on $\thetamodel$ (the dominating measure), the log-likelihood
is
\begin{equation}\label{eq:ch03:jacod}
\loglik_{\Tend}(\thetamodel)
 \;=\;
 \sum_{m=1}^{n} \log \lam_{e_m}(t_m;\thetamodel)
 \;-\;
 \sum_{e\in\mathcal E} \int_0^{\Tend} \atrisk_e(u)\,\lam_e(u;\thetamodel)\dd u ,
\end{equation}
which is \cref{eq:ch02:jacod} of \cref{def:ch02:intensity}.
\end{theorem}

\begin{proof}
We give the product-integral derivation over infinitesimal bins; it is the
argument to remember, and it becomes rigorous as a product-integral limit
\citep{daley2008,abgk1993}. Partition $[0,\Tend]$ into $K$ bins
$B_k = ((k-1)\delta, k\delta]$ with $\delta = \Tend/K$, and $K$ large
enough that no bin contains two events. By the chain rule of conditional
probability, the probability of the observed configuration is the product
over bins of the conditional probability of what happened in $B_k$ given
$\F_{(k-1)\delta}$. By \cref{eq:ch03:heuristic} and
Assumption~\ref{asm:ch03:regularity}, a bin containing the event $(t_m, e_m)$
contributes $\lam_{e_m}(t_m)\,\delta + o(\delta)$; an empty bin contributes
$1 - \sum_e \atrisk_e\lam_e\,\delta + o(\delta)$; two or more events in one
bin have probability $o(\delta)$. Hence
\[
\log \Prob(\text{data}) - n\log\delta
 \;=\;
 \sum_{m=1}^{n} \log\lam_{e_m}(t_m)
 \;+\;
 \sum_{k:\,B_k \text{ empty}} \log\Big(1 - \textstyle\sum_e \atrisk_e\lam_e\,\delta\Big)
 \;+\; o(1).
\]
Since $\log(1-x) = -x + O(x^2)$, the middle sum is a Riemann sum of
$-\sum_e \int_0^{\Tend}\atrisk_e\lam_e\dd u$ and converges to it as
$\delta \to 0$. The diverging term $n\log\delta$ is the same for every
$\thetamodel$: it cancels in every likelihood ratio and belongs to the
dominating measure (concretely, \cref{eq:ch03:jacod} exponentiated and
multiplied by $e^{\Tend}$ is the density of the observed process with
respect to a unit-rate Poisson process). This yields
\cref{eq:ch03:jacod}.
\end{proof}

The two terms have distinct jobs. The \emph{event term}
$\sum_m \log\lam_{e_m}(t_m)$ rewards the model for placing intensity where
events actually occurred; it is evaluated exactly at the jumps, on the log
scale, so assigning near-zero intensity to a realized event is very
costly. The \emph{compensator term} $\sum_e\int \atrisk_e\lam_e$ charges
the model for exposure: every unit of intensity claimed over every
firm-week at risk is paid for linearly, whether or not an event arrives. A
model can only win the event term by spending in the compensator term, and
the optimal spending pattern is the true intensity --- the content of the
score identity in \cref{ch03:sec:estimation}. Two consequences of the
shape of \cref{eq:ch03:jacod} recur constantly. An event-free spell
$(a,b]$ contributes exactly the survival factor
$\exp(-\int_a^b \atrisk\lam)$; censored spells therefore carry information
and must never be dropped (\cref{warn:ch02:openspells},
\cref{prop:ch03:censoring}). And the event term is the only place event
\emph{times} enter; if the times are corrupted, the compensator term
survives while the event term degrades --- the starting point for the
interval-censored treatment of \cref{ch:censoring}.

\section{The ground/mark factorization}\label{ch03:sec:groundmark}

The market process of \cref{def:ch02:market} is a \emph{marked} point
process: ground times $t_m$ with marks $\dmark_m$ (in our data the mark is
the triple sector--firm--terms). The likelihood of a marked process
factorizes exactly into a timing part and a mark part. This identity is
the formal license for the two-layer architecture, so we set it up
carefully and prove it.

\begin{definition}[Marked intensity kernel]\label{def:ch03:markedintensity}
Let marks live in a measurable space $(\mathcal M, \nu)$ with $\nu$ a
$\sigma$-finite reference measure (counting measure for discrete marks such
as sector and firm identity, Lebesgue for continuous terms). A
\emph{marked point process} with intensity kernel $\lam(t,\dmark)$ is one
for which, for every measurable $A \subseteq \mathcal M$, the counting
process $N(t,A) = \#\{m : t_m \le t,\ \dmark_m \in A\}$ has conditional
intensity $\int_A \lam(t,\dmark)\,\nu(\dd\dmark)$. The \emph{ground
process} is $N_{\mathrm g}(t) = N(t,\mathcal M)$ with \emph{ground
intensity} $\lam_{\mathrm g}(t) = \int_{\mathcal M}
\lam(t,\dmark)\,\nu(\dd\dmark)$, and the \emph{mark density} is
\begin{equation}\label{eq:ch03:markdensity}
p\big(\dmark \st t, \Hist_{t^-}\big)
 \;=\; \frac{\lam(t,\dmark)}{\lam_{\mathrm g}(t)}
 \qquad \text{whenever } \lam_{\mathrm g}(t) > 0,
\end{equation}
a probability density in $\dmark$ with respect to $\nu$, predictable in
$t$.
\end{definition}

\begin{theorem}[Ground/mark factorization]\label{thm:ch03:groundmark}
Under Assumption~\ref{asm:ch03:regularity} applied to the family
$\{N(\cdot,A)\}$, the log-likelihood of the marked data
$\{(t_m,\dmark_m)\}_{m=1}^{n}$ on $[0,\Tend]$ satisfies, exactly and for
every parameter value and every realization,
\begin{equation}\label{eq:ch03:groundmark}
\loglik_{\Tend}
 \;=\;
 \underbrace{\;\sum_{m=1}^{n} \log \lam_{\mathrm g}(t_m)
 \;-\; \int_0^{\Tend} \lam_{\mathrm g}(u)\dd u\;}_{\text{ground factor: timing}}
 \;+\;
 \underbrace{\;\sum_{m=1}^{n} \log p\big(\dmark_m \st t_m, \Hist_{t_m^-}\big)\;}_{\text{mark factor: identity}} .
\end{equation}
\end{theorem}

\begin{proof}
Take first the case of finitely many mark values, which covers the
sector--firm part of our data. The marked process is then the multivariate
counting process $(N(\cdot,\{\dmark\}))_{\dmark\in\mathcal M}$ with
intensities $\lam(t,\dmark)$, and \cref{thm:ch03:jacod} gives
\[
\loglik_{\Tend}
 \;=\;
 \sum_{m=1}^{n} \log \lam(t_m, \dmark_m)
 \;-\;
 \sum_{\dmark\in\mathcal M} \int_0^{\Tend} \lam(u,\dmark)\dd u .
\]
Substitute $\lam(t,\dmark) = \lam_{\mathrm g}(t)\,
p(\dmark\st t,\Hist_{t^-})$ from \cref{eq:ch03:markdensity}. The event
term splits into $\sum_m \log\lam_{\mathrm g}(t_m) + \sum_m \log
p(\dmark_m \st t_m, \Hist_{t_m^-})$. The compensator term collapses:
\[
\sum_{\dmark} \int_0^{\Tend} \lam_{\mathrm g}(u)\,
 p(\dmark \st u, \Hist_{u^-})\dd u
 \;=\;
 \int_0^{\Tend} \lam_{\mathrm g}(u)
 \underbrace{\sum_{\dmark} p(\dmark \st u, \Hist_{u^-})}_{=\,1}\dd u
 \;=\;
 \int_0^{\Tend} \lam_{\mathrm g}(u)\dd u ,
\]
because the mark density integrates to one at every $u$. For general mark
spaces, replace the sum over $\dmark$ by $\int \cdot\,\nu(\dd\dmark)$; the
computation is identical and the underlying likelihood formula is Jacod's
formula for marked processes \citep{daley2008}. This proves
\cref{eq:ch03:groundmark}.
\end{proof}

The crux is worth stating in words: \emph{the mark density leaves no trace
in the compensator}, because it integrates to one under the exposure
integral. The mark factor therefore contributes only at event times --- it
is a product of genuine conditional probabilities ``given a deal at $t_m$,
which mark'' --- and involves intensities only through the ratio
\cref{eq:ch03:markdensity}, in which any common timing factor cancels.
That is the mathematical root of the robustness asymmetry announced in
\cref{ch01:sec:factorization}: distortions that rescale or warp
$\lam_{\mathrm g}$ (untrusted dates, aggregation, baseline
misspecification) damage the ground factor and cancel from the mark
factor. \Cref{ch:censoring,ch:ranking} quantify this.

\begin{remark}[Identity, not approximation]\label{rem:ch03:identity}
\Cref{eq:ch03:groundmark} assumes no independence, no Poisson structure,
and no correct specification; it is algebra performed on the likelihood of
\cref{thm:ch03:jacod}, valid realization by realization. Every
approximation made in this document lives \emph{inside} one factor: the
weekly MBPP approximation of \cref{ch:censoring} approximates the ground
factor; the linear-index softmax of \cref{ch:ranking} approximates the
mark factor; neither touches the other factor, and a data defect can be
assigned to the factor whose term it corrupts. This bookkeeping --- which
factor is damaged --- is how the defect ledger of \cref{tab:ch02:defects}
becomes actionable.
\end{remark}

\begin{corollary}[The two-layer sector$\to$firm architecture]\label{cor:ch03:twolayer}
Specialize the mark to $\dmark = (s, i, \dmark')$: sector, receiving firm,
deal terms. Applying \cref{thm:ch03:groundmark} and factoring the mark
density by the chain rule,
$p(s,i,\dmark' \st t) = p(s \st t)\, p(i \st s, t)\, p(\dmark' \st s,i,t)$,
the likelihood splits into the three factors of \cref{eq:ch01:groundmark}.
Equivalently, defining sector processes $N_s$ with intensities
$\lam_s(t) = \lam_{\mathrm g}(t)\,p(s\st t,\Hist_{t^-})$, every tracked
firm's intensity factorizes as
\begin{equation}\label{eq:ch03:twolayer}
\lam_i(t)
 \;=\;
 \lam_{s(i)}(t)\;
 \Prob\big(i \;\big|\; s(i),\, t,\, \Hist_{t^-}\big),
 \qquad i \in \firms,
\end{equation}
and the log-likelihood is the sum of a sector-timing part (in
$\lam_s$ only) and a firm-identity part (in the conditional choice
probabilities only).
\end{corollary}

\begin{proof}
Apply the theorem twice: once with mark $(s,i,\dmark')$, once viewing $i$
as a mark of the sector-$s$ ground process $N_s$; the chain-rule
factorization of the joint mark density is the elementary identity for
conditional densities, and \cref{eq:ch03:twolayer} is
\cref{eq:ch03:markdensity} read at the firm level.
\end{proof}

\begin{decision}{D3.1 --- the factorization is the architecture}
The repository's two-layer design --- a sector-level ground model
(\modrepo{sector\_stability.py}, \modrepo{eventtime/}) and a firm-level
risk-set choice model (\modrepo{sector\_ranker.py}) --- is
\cref{cor:ch03:twolayer} implemented literally, with the two factors
parameterized by variation-independent parameter blocks so that maximizing
the sum \cref{eq:ch03:groundmark} decouples into two exact, separate
maximizations. Rationale: the factorization is an identity, the factors'
metrics align with distinct targets (T2 vs.\ T1 of
\cref{ch01:sec:targets}), and the factors have opposite robustness
profiles under the defect ledger. Reversal trigger: evidence of material
parameter sharing across factors --- e.g.\ a jointly estimated timing--identity
model beating the separate fits on \emph{both} held-out ground NLL and
identity recall@$K$ in two consecutive evaluation campaigns --- would move
us to joint estimation; the factorized \emph{likelihood} would remain, only
the separability of its maximization would be given up.
\end{decision}

\section{Superposition and the identity of the next event}\label{ch03:sec:superposition}

Two small results connect the multivariate and marked views, and both are
used constantly --- in the ranking chapter as the modeling target, and in
the simulation section as the attribution step.

\begin{proposition}[Superposition]\label{prop:ch03:superposition}
Under Assumption~\ref{asm:ch03:regularity}, the pooled process
$\bar N = \sum_{e} N_e$ is a counting process with conditional intensity
$\bar\lam(t) = \sum_e \atrisk_e(t)\lam_e(t)$.
\end{proposition}

\begin{proof}
$\bar N$ is right-continuous, nondecreasing, and has unit jumps because no
two components jump together; and $\bar N - \int \bar\lam
= \sum_e M_e$ is a finite sum of local martingales, hence a local
martingale. Since $\bar\lam$ is a sum of predictable processes it is
predictable, so it is the conditional intensity of $\bar N$.
\end{proof}

\begin{proposition}[Which type fires]\label{prop:ch03:nexttype}
Under Assumption~\ref{asm:ch03:regularity}, conditionally on $\Ftm$ and on the event
that the pooled process jumps at $t$,
\begin{equation}\label{eq:ch03:nexttype}
\Prob\big(\text{the jump is of type } e \;\big|\; \Ftm,\ \Delta\bar N(t)=1\big)
 \;=\;
 \frac{\atrisk_e(t)\,\lam_e(t)}{\sum_{e'} \atrisk_{e'}(t)\,\lam_{e'}(t)} .
\end{equation}
\end{proposition}

\begin{proof}[Proof sketch]
By \cref{eq:ch03:heuristic}, $\Prob(\text{type-}e \text{ jump in }
[t,t+\dd t) \st \Ftm) = \atrisk_e\lam_e(t)\dd t$ and
$\Prob(\text{some jump}) = \bar\lam(t)\dd t$; the ratio gives
\cref{eq:ch03:nexttype} as $\dd t \to 0$. Structurally, the statement is
\cref{thm:ch03:groundmark} read backwards: view the type $e$ as the mark
of the pooled process; then the right-hand side of
\cref{eq:ch03:nexttype} is exactly the mark density
\cref{eq:ch03:markdensity}. The full measure-theoretic proof, together
with its operational consequences --- the risk-set conditional logit, its
partial-likelihood justification \citep{cox1975,mcfadden1974}, and the
evaluation protocol built on it --- is developed in \cref{ch:ranking},
which restates \cref{eq:ch03:nexttype} as the definition of the ranking
target.
\end{proof}

\Cref{eq:ch03:nexttype} is the entire T1 target of
\cref{ch01:sec:targets} in one line: ranking firms by
$\lam_i(t)/\sum_j \lam_j(t)$ is ranking by conditional next-event
probability, and by \cref{eq:ch03:twolayer} the sector timing factor
cancels from the ratio within a sector. It is also the attribution rule of
the thinning simulator below: accept a proposed event, then label it type
$e$ with probability $\lam_e/\bar\lam$.

\section{Martingale estimation theory}\label{ch03:sec:estimation}

Why does maximizing \cref{eq:ch03:jacod} work, given that events are
dependent, non-stationary, and censored? The answer is not the classical
i.i.d.\ theory but the martingale structure of
\cref{eq:ch03:compensator}; this section develops exactly the pieces later
chapters invoke.

Differentiating \cref{eq:ch03:jacod} in $\thetamodel$ (legitimate under
Assumption~\ref{asm:ch03:regularity}(iii)) and writing
$\dd N_e - \atrisk_e \lam_e \dd u = \dd M_e$ gives the \emph{score
process}
\begin{equation}\label{eq:ch03:score}
\scorefn(\thetamodel, t)
 \;=\;
 \sum_{e} \int_0^{t} \partial_{\thetamodel} \log \lam_e(u;\thetamodel)
 \,\big[\dd N_e(u) - \atrisk_e(u)\lam_e(u;\thetamodel)\dd u\big]
 \;=\;
 \sum_{e} \int_0^{t} \partial_{\thetamodel} \log \lam_e \,\dd M_e(\thetamodel),
\end{equation}
where the second equality holds when $\thetamodel$ is the true value: the
score is a stochastic integral of a predictable process against the
compensated noise.

\begin{proposition}[Unbiased score at the truth]\label{prop:ch03:score}
Let $\thetamodel_0$ be the true parameter. Under
Assumption~\ref{asm:ch03:regularity}, $\scorefn(\thetamodel_0,\cdot)$ is a mean-zero
local square-integrable martingale; under the integrability of
Assumption~\ref{asm:ch03:regularity}(ii)--(iii) it is a martingale, and in
particular
$\E\,\scorefn(\thetamodel_0,\Tend) = 0$.
\end{proposition}

\begin{proof}
At $\thetamodel_0$, each $M_e$ in \cref{eq:ch03:score} is a local
martingale by \cref{eq:ch03:compensator}, and the integrand
$H_e = \partial_{\thetamodel}\log\lam_e$ is predictable and locally
bounded by assumption. A stochastic integral of a predictable locally
bounded process against a local martingale is a local martingale
\citep[Ch.~II]{abgk1993}, and $\scorefn(\thetamodel_0,0)=0$. The moment
condition upgrades the local martingale to a martingale, whence the
expectation is zero at every $t$, in particular at $\Tend$. What breaks
when the assumptions fail is instructive: if $\lam$ peeks at time-$t$
information (\cref{warn:ch03:lookahead}), $H_e$ is not predictable and the
integral is no longer a martingale --- the score acquires a bias of the
same sign as the leak; if the model is misspecified, $M_e(\thetamodel)$ is
not a martingale for any $\thetamodel$ in the family, and the score is
instead centered at the pseudo-true parameter $\thetamodel_*$ --- the
Kullback--Leibler projection of the truth onto the family --- which is
what the sandwich variance below is built for.
\end{proof}

\begin{proposition}[Information identities]\label{prop:ch03:information}
Define the observed information
$J(\thetamodel) = -\partial^2_{\thetamodel}\loglik_{\Tend}(\thetamodel)$
and the predictable variation of the score at the truth,
\begin{equation}\label{eq:ch03:predvar}
\langle \scorefn(\thetamodel_0,\cdot)\rangle(\Tend)
 \;=\;
 \sum_{e}\int_0^{\Tend}
 \big(\partial_{\thetamodel}\log\lam_e\big)
 \big(\partial_{\thetamodel}\log\lam_e\big)\T
 \,\atrisk_e\,\lam_e \dd u .
\end{equation}
Then $J(\thetamodel_0) - \langle
\scorefn(\thetamodel_0,\cdot)\rangle(\Tend) = -\sum_e \int_0^{\Tend}
\partial^2_{\thetamodel}\log\lam_e \,\dd M_e$ is a mean-zero martingale
term, and consequently (the second Bartlett identity)
$\E\,J(\thetamodel_0) = \E\,\langle\scorefn\rangle(\Tend)
 \eqqcolon \Fisher(\thetamodel_0)$,
the expected Fisher information.
\end{proposition}

\begin{proof}
Differentiate \cref{eq:ch03:jacod} twice:
$-J = \sum_e[\int \partial^2\log\lam_e \dd N_e - \int
\atrisk_e\,\partial^2\lam_e \dd u]$. Substitute
$\partial^2\lam_e = \lam_e\{\partial^2\log\lam_e +
(\partial\log\lam_e)(\partial\log\lam_e)\T\}$ in the second integral and
collect terms:
$J = -\sum_e\int \partial^2\log\lam_e\,\dd M_e +
\langle\scorefn\rangle(\Tend)$. The first term is a stochastic integral of
a predictable process against a martingale, mean zero at
$\thetamodel_0$ by the argument of \cref{prop:ch03:score}; take
expectations. That the compensator of the score's quadratic variation is
\cref{eq:ch03:predvar} is the standard predictable-variation computation
for counting-process integrals \citep[Ch.~II]{abgk1993}.
\end{proof}

\paragraph{Sandwich variance.}
The MLE $\thetahat$ solves $\scorefn(\thetahat,\Tend)=0$; a Taylor
expansion around the (pseudo-)truth gives $\thetahat - \thetamodel_*
\approx J(\thetamodel_*)^{-1}\scorefn(\thetamodel_*,\Tend)$, hence the
asymptotic variance
\begin{equation}\label{eq:ch03:sandwich}
\Var(\thetahat) \;\approx\; A^{-1} B\, A^{-1},
\qquad
A = \E\,J(\thetamodel_*),
\quad
B = \Var\,\scorefn(\thetamodel_*,\Tend).
\end{equation}
When the model is correctly specified, \cref{prop:ch03:information} gives
$B = A = \Fisher(\thetamodel_0)$ and \cref{eq:ch03:sandwich} collapses to
the familiar $\Fisher^{-1}$. When it is not --- and we treat
misspecification as the operating condition; regime B of the synthetic
world (\cref{ch02:sec:synthetic}) is misspecified by construction --- the
equality fails and only the sandwich is valid. In practice $B$ is
estimated by the optional-variation analogue of \cref{eq:ch03:predvar}
(outer products of per-event score contributions), clustered at the sector
level because latent frailty is expected to correlate firms within a
sector; the frailty experiments planned in \cref{ch:results-synth} will
measure how much this matters.

\begin{theorem}[Rebolledo's martingale CLT; statement only]\label{thm:ch03:rebolledo}
Let $\scorefn^{(n)}$ be a sequence of mean-zero local square-integrable
martingales. If (i) $\langle \scorefn^{(n)}\rangle(t) \to V(t)$ in
probability for each $t$, with $V$ deterministic and continuous, and (ii)
for each $\eps>0$ the predictable variation contributed by jumps of
magnitude exceeding $\eps$ vanishes in probability, then
$\scorefn^{(n)}$ converges weakly to a continuous Gaussian martingale with
variance function $V$ \citep{abgk1993,andersengill1982}.
\end{theorem}

Rebolledo's theorem is the replacement for the classical CLT in all
counting-process asymptotics: applied to the normalized score with
$V = $ the limit of \cref{eq:ch03:predvar}, and combined with the Taylor
expansion above, it yields asymptotic normality of $\thetahat$ with the
variance \cref{eq:ch03:sandwich}. Condition (ii) holds here because
individual events contribute vanishing shares of the total information as
the window or the cross-section grows.

\begin{theorem}[Consistency and asymptotic normality of the point-process MLE; statement only]\label{thm:ch03:mle}
For stationary ergodic point processes with intensity families satisfying
identifiability, smoothness, and moment conditions,
$\thetahat_{\Tend} \to \thetamodel_0$ a.s.\ and
$\sqrt{\Tend}\,(\thetahat_{\Tend} - \thetamodel_0) \Rightarrow
\mathcal N\big(0, \Fisher_1(\thetamodel_0)^{-1}\big)$ as
$\Tend \to \infty$, with $\Fisher_1$ the per-unit-time information
\citep{ogata1978}. The counting-process regime with many units observed on
a fixed window --- our firm panel --- is covered by the same martingale
program \citep{andersengill1982,abgk1993}.
\end{theorem}

Both asymptotic regimes are present in this repository and should not be
conflated: sector-level timing models see $S = 12$ fixed components over a
long window (Ogata's $\Tend\to\infty$ regime), while the mark factor sees
thousands of firms per risk set on a fixed window (the Andersen--Gill
$n\to\infty$ regime). Standard errors quoted in later chapters state which
regime they invoke.

\begin{decision}{D3.2 --- martingale-grade inference outputs}
Every likelihood fitter in \modrepo{} must expose (i) the analytic score
\cref{eq:ch03:score}, (ii) an observed-information or numerical Hessian,
and (iii) sandwich standard errors \cref{eq:ch03:sandwich} with
sector-level clustering as the default reported uncertainty; model-based
$\Fisher^{-1}$ errors may be shown only alongside. Rationale:
misspecification is the operating condition (regime B), under which only
the sandwich is valid for the pseudo-true parameter; and exposing the
score is what makes the gradient checks and the GOF machinery of
\cref{ch:gof} cheap. Reversal trigger: for a model family that the
diagnostics of \cref{ch:gof} certify as well specified on real data across
two campaigns, model-based errors may be promoted to primary for that
family only.
\end{decision}

\section{The time-rescaling theorem}\label{ch03:sec:rescaling}

If the fitted intensity is right, rescaling time by the fitted compensator
turns the data into a unit-rate Poisson process. This is the
goodness-of-fit engine of \cref{ch:gof}; we state and prove it here.

\begin{theorem}[Time rescaling]\label{thm:ch03:rescaling}
Let $N$ have conditional intensity $\lam$ under
Assumption~\ref{asm:ch03:regularity}, with compensator
$\Lam(t) = \int_0^t \atrisk\lam\dd u$ continuous and strictly increasing
on the observation window, and let $t_1 < t_2 < \cdots < t_n$ be the event
times on $[0,\Tend]$. Then the rescaled times $\tau_m = \Lam(t_m)$ form a
unit-rate Poisson process on $[0, \Lam(\Tend)]$; equivalently, the
rescaled gaps
\begin{equation}\label{eq:ch03:rescaled}
\Delta_m \;=\; \Lam(t_m) - \Lam(t_{m-1})
 \;=\; \int_{t_{m-1}}^{t_m} \atrisk(u)\lam(u)\dd u,
\qquad t_0 = 0,
\end{equation}
are i.i.d.\ $\mathrm{Exp}(1)$, and $\Lam(\Tend) - \Lam(t_n)$ is an
independent right-censored $\mathrm{Exp}(1)$ variable
\citep{brown2002,ogata1988}.
\end{theorem}

\begin{proof}
We give the hazard argument in the case where, between events, the
intensity path is determined by the history at the last event together
with the observed covariate paths --- true of every model fitted in this
document, whose $\lam(u)$ for $u \in (t_{m-1}, t_m]$ is an explicit
functional of $\Hist_{t_{m-1}}$, the no-further-event continuation, and
covariates that $\F$ carries. Condition on $\F_{t_{m-1}}$ and let $g(u)$
denote that determined intensity path. By \cref{eq:ch03:heuristic}, the
waiting time to the next event has hazard $g$:
$\Prob(t_m \in [u, u+\dd u) \st t_m \ge u, \F_{t_{m-1}}) = g(u)\dd u$,
hence the conditional survivor function
$\Prob(t_m > u \st \F_{t_{m-1}}) = \exp(-\int_{t_{m-1}}^{u} g)$. Now fix
$x \ge 0$ and let $u_x = \inf\{u : \int_{t_{m-1}}^{u} g = x\}$, finite and
uniquely attained because $\Lam$ is continuous and strictly increasing.
The events $\{\Delta_m > x\}$ and $\{t_m > u_x\}$ coincide, so
\[
\Prob\big(\Delta_m > x \;\big|\; \F_{t_{m-1}}\big)
 \;=\;
 \exp\Big(-\int_{t_{m-1}}^{u_x} g(u)\dd u\Big)
 \;=\; e^{-x}.
\]
The conditional law of $\Delta_m$ given the entire past is $\mathrm{Exp}(1)$,
free of the past; iterating over $m$ makes $(\Delta_1,\dots,\Delta_n)$
i.i.d.\ $\mathrm{Exp}(1)$, which is the gap characterization of the
unit-rate Poisson process. The same computation with $x$ at the censoring
boundary gives the final censored gap. In full generality --- intensity
paths influenced by additional external randomness between events --- the
result is the random-time-change theorem: $M(\Lam^{-1}(\tau))$ is a
martingale by optional stopping, so the time-changed counting process has
deterministic continuous compensator $\tau$, and Watanabe's
characterization identifies it as unit-rate Poisson \citep{daley2008}.
\end{proof}

The practical reading: \emph{a correct model makes its own residuals
exponential}, whatever the dependence structure of the original data. The
resulting protocol --- rescale with the fitted compensator, then test the
gaps against $\mathrm{Exp}(1)$ by Kolmogorov--Smirnov and QQ plots, as in
Ogata's residual analysis \citep{ogata1988} --- is fixed as this
repository's primary event-time diagnostic in \cref{ch:gof}, including
the treatment of the censored last gap. Two caveats travel with it. The
test consumes exact event times, so under defect (C4) --- weekly
aggregation --- it degrades by construction; the interval-censored regime
uses bucket-level Pearson and dispersion residuals instead
(\cref{ch:gof}, \modrepo{mbpp/gof.py}). And rescaling with an intensity
fitted on the \emph{same} data is optimistic; the split protocol of
\cref{ch:gof} rescales held-out windows with frozen parameters.

\section{Simulation: inversion and thinning}\label{ch03:sec:simulation}

Simulation matters here because the synthetic world of
\cref{ch02:sec:synthetic}, the parameter-recovery studies planned in
\cref{ch:results-synth}, and the posterior predictive checks of
\cref{ch:gof} all rest on generators whose output provably has the
intended intensity. Two algorithms cover our needs.

\begin{corollary}[Inverse-compensator simulation]\label{cor:ch03:inversion}
Let $E_1, E_2, \dots$ be i.i.d.\ $\mathrm{Exp}(1)$. Setting, recursively,
\[
t_m \;=\; \inf\Big\{u > t_{m-1} : \int_{t_{m-1}}^{u} \lam(v)\dd v \ge E_m\Big\},
\]
with $\lam$ evaluated along the history generated so far, produces a point
process with conditional intensity $\lam$.
\end{corollary}

\begin{proof}
This is \cref{thm:ch03:rescaling} read backwards: the construction forces
the rescaled gaps to be the i.i.d.\ $\mathrm{Exp}(1)$ inputs, and the
hazard computation in that proof, run in reverse, shows the waiting time
generated has exactly the conditional hazard $\lam$.
\end{proof}

Inversion is exact and fast whenever the conditional compensator inverts
in closed form --- the covariate-driven Poisson models of
\cref{ch:poisson} with piecewise-constant covariates are the leading case.
When the compensator does not invert cheaply, we thin.

\begin{proposition}[Correctness of Ogata's modified thinning]\label{prop:ch03:thinning}
Suppose that on each inter-refresh interval (between consecutive accepted
events, covariate breakpoints, or bound refreshes) a constant
$\bar\lam$ is available with $\lam(u) \le \bar\lam$ for all $u$ in the
interval, where $\lam$ is evaluated along the accepted history and its
no-new-event continuation. Generate proposals on the interval from a
homogeneous Poisson process of rate $\bar\lam$, and accept a proposal at
$\tau$ independently with probability $\lam(\tau)/\bar\lam$, where
$\lam(\tau)$ is computed from the accepted history up to $\tau^-$. Then
the accepted points form a point process with conditional intensity
$\lam$ \citep{lewisshedler1979,ogata1981}.
\end{proposition}

\begin{proof}
Work on the enlarged filtration generated by the accepted history, the
proposal process, and the acceptance variables. Given the accepted history
$\F^{\mathrm{acc}}_{t^-}$, the probability of an accepted point in
$[t, t+\dd t)$ is the probability of a proposal there,
$\bar\lam\dd t + o(\dd t)$, times the acceptance probability
$\lam(t)/\bar\lam$ --- which is $\F^{\mathrm{acc}}_{t^-}$-measurable
(predictable), because $\lam(t)$ depends only on accepted points before
$t$ and on covariates; two proposals in the bin have probability
$o(\dd t)$. The product is $\lam(t)\dd t + o(\dd t)$, so the accepted
counting process minus $\int \lam$ is a martingale on the enlarged
filtration, and since $\lam$ is adapted to the accepted history, it is
the conditional intensity of the accepted process
(\cref{eq:ch03:compensator}). The bound being valid on the whole
inter-refresh interval is what makes the acceptance probability well
defined ($\le 1$) throughout; the multivariate version accepts with
$\bar\lam^{-1}\sum_e \lam_e(\tau)$ and then labels the accepted point type
$e$ with probability $\lam_e(\tau)/\sum_{e'}\lam_{e'}(\tau)$, which is
correct by \cref{prop:ch03:nexttype}.
\end{proof}

\begin{warning}[The bound is part of the algorithm's correctness]\label{warn:ch03:bound}
Thinning is exact only while $\bar\lam$ truly dominates. For exponential
Hawkes with nonnegative kernels ($\kernel(t) = \branch\decay
e^{-\decay t}$), the intensity strictly decreases between events, so the
value just after the last event is a valid bound --- \emph{unless the
baseline jumps upward inside the interval}, as it does whenever a
piecewise-constant covariate steps. A generator that fails to refresh the
bound at covariate breakpoints silently under-simulates events after
every upward step, biasing any simulation-based conclusion toward weaker
covariate effects. \modrepo{eventtime/simulate.py} refreshes at
accepted events \emph{and} at covariate breakpoints for exactly this
reason, and verifies its recursive intensity against a naive
recomputation.
\end{warning}

For the exponential kernel there is also an exact, thinning-free
simulator that exploits the Markov property of the intensity itself
(cluster and shot-noise structure), due to \citet{dassioszhao2013}; it and
the branching (cluster) representation are developed where they are used,
in \cref{ch:hawkes}.

\section{What the machinery absorbs, and what it cannot}\label{ch03:sec:defects}

Decision D2.1 of \cref{ch:data} requires every model chapter to declare
its position against the defect ledger of \cref{tab:ch02:defects}. This
chapter contributes the declaration for the machinery itself: which
defects the counting-process framework absorbs \emph{by construction}, and
which no amount of likelihood theory can repair.

\begin{proposition}[Independent censoring preserves the likelihood; (C1)]\label{prop:ch03:censoring}
Let $C_e$ be $\{0,1\}$-valued predictable censoring processes, and let
$\mathcal G_t \supseteq \Ft$ be the filtration enlarged with the censoring
information. Assume independent censoring: each $N_e$ has the same
intensity $\atrisk_e\lam_e$ with respect to $\{\mathcal G_t\}$ as with
respect to $\{\Ft\}$. Then the observed processes
$N_e^{c}(t) = \int_0^t C_e \dd N_e$ have $\mathcal G$-intensity
$C_e\,\atrisk_e\,\lam_e$, and \cref{eq:ch03:jacod} computed on the
observed events with exposure $C_e\atrisk_e$ is a genuine (partial)
likelihood: its score retains the martingale property of
\cref{prop:ch03:score}, and all of \cref{ch03:sec:estimation} applies
unchanged \citep[III.2]{abgk1993}.
\end{proposition}

\begin{proof}[Proof sketch]
Since $C_e$ is predictable, its value at $t$ is known given
$\mathcal G_{t^-}$, so
$\E[\dd N^c_e \st \mathcal G_{t^-}] = C_e(t)\,\E[\dd N_e \st \mathcal
G_{t^-}] = C_e\,\atrisk_e\,\lam_e\dd t$, the middle equality being the
independent-censoring hypothesis. Hence
$N^c_e - \int C_e \atrisk_e \lam_e$ is a $\mathcal G$-martingale, and
\cref{thm:ch03:jacod} applied to $(N^c_e)$ yields
\cref{eq:ch03:jacod} with $\atrisk_e$ replaced by $C_e\atrisk_e$; at
observed events $C_e = 1$, so the event term is untouched. The factor of
the joint likelihood governing the censoring mechanism itself contains no
$\thetamodel$ and is dropped --- that is the sense in which the result is
a \emph{partial} likelihood. Administrative censoring at fixed $\Tend$
has $C_e(t) = \ind{t \le \Tend}$ deterministic, hence predictable, and
trivially preserves the intensity: it is independent by construction.
What breaks without predictability is the first equality; what breaks
without independence is the second --- if censoring anticipates the
process (firms leave observation \emph{because} their funding prospects
changed), $\E[\dd N_e \st \mathcal G_{t^-}] \neq \atrisk_e\lam_e\dd t$ and
the observed-data likelihood is biased. That failure mode is exactly
defect (C3).
\end{proof}

Operationally, \cref{prop:ch03:censoring} is the license already
exercised in \cref{warn:ch02:openspells}: a spell open at June 2024
contributes its survival factor $\exp(-\int \lam)$ through the
compensator term, no term is dropped and none invented, \emph{provided
the open spell is kept in the data}. The machinery is robust to (C1).

\paragraph{(C2) Left truncation.}
Delayed entry is absorbed by the at-risk indicator: setting
$\atrisk_i(t) = \ind{\entry_i \le t < U_i}$ makes every compensator
integral in \cref{eq:ch03:jacod} start at the entry time $\entry_i$, and
the likelihood is valid \emph{conditionally on the pre-entry history}
\citep{keiding1990,kalbfleisch2002}. The mechanical fix is complete; the
statistical selection is not fixed at all: because entry occurs
\emph{by raising}, every estimand is conditional on ``at least one deal
before $\Tend$,'' and unconditional statements about companies require
the open-universe treatment of \cref{ch:universe}. The machinery
\emph{corrects (C2) under the stated conditional interpretation}, and
silently changes the estimand if that interpretation is forgotten.

\paragraph{(C3) Unlabeled exits.}
Here the machinery is \emph{vulnerable, with a known sign}. A dead firm
with $\atrisk_i \equiv 1$ contributes exposure and no events ---
indistinguishable, inside the likelihood, from a live quiet firm --- so no
score, residual, or information identity can detect the corruption:
\cref{eq:ch03:jacod} is being evaluated with the wrong $\atrisk$, and
every fitted intensity is biased \emph{downward} while ranking risk sets
are contaminated with phantom competitors. We expect the damage to be
large; the synthetic study of \cref{ch:results-synth} is designed to
measure it.
Treatments are exposure windows and mover--stayer structure in
\cref{ch:poisson,ch:buckets} and ranking-layer mitigations in
\cref{ch:ranking}; none is a consequence of the machinery, all are
assumptions imported to compensate for it.

\paragraph{(C4) Untrusted times and aggregation.}
The event term of \cref{eq:ch03:jacod} --- and only the event term ---
consumes exact times, and \cref{thm:ch03:rescaling} and
\cref{prop:ch03:thinning} consume them too. Date jitter perturbs
$\log\lam(t_m)$ where $\lam$ moves fastest, i.e.\ it attacks timescale
parameters first; full aggregation deletes the event term and requires
the interval-censored likelihoods of \cref{ch:censoring}
\citep{rizoiu2022}. Our working hypothesis, to be tested in
\cref{ch:results-synth}, is that $\branch$ survives aggregation while
$\decay$ does not. The machinery is vulnerable to (C4);
\cref{ch:censoring} is the correction.

\paragraph{(C5) and (C6).}
Stale covariates leave the machinery \emph{valid but reinterpreted}: LOCF
covariates as constructed under \cref{prot:ch02:pit} are predictable, so
\cref{eq:ch03:jacod} is a genuine likelihood for the model that
conditions on the \emph{observed} covariates --- but its coefficients are
attenuated relative to the true-covariate model, which is
\cref{ch:stale}'s subject. Out-of-universe deals leave sector-level
ground counts complete but thin the firm-level mark factor: without an
explicit outside option, the normalization in
\cref{eq:ch03:nexttype} overstates every tracked firm's probability ---
\cref{ch:universe} owns the fix.

In summary: the machinery of this chapter is robust to (C1), corrects
(C2) under a conditional reading, and is structurally blind to (C3),
(C4) on the event term, (C5) in interpretation, and (C6) in the mark
normalization. That blindness is the agenda of
\cref{ch:poisson,ch:censoring,ch:stale,ch:universe}.

\begin{codehooks}
\Cref{thm:ch03:jacod} is implemented in \modrepo{eventtime/likelihood.py}:
\code{log\_likelihood} assembles the event term by recursive
exponential-kernel sums decayed to each event's left limit
(\cref{warn:ch03:lookahead}) plus a closed-form compensator term;
\code{log\_likelihood\_and\_grad} returns the analytic score
\cref{eq:ch03:score}, satisfying decision D3.2.
\Cref{prop:ch03:thinning,warn:ch03:bound} are implemented in
\modrepo{eventtime/simulate.py} (\code{simulate\_multivariate\_hawkes}):
Ogata thinning with bound refresh at accepted events and covariate
breakpoints, with a naive-recomputation cross-check of the recursive
intensity. \Cref{thm:ch03:rescaling} lives in \modrepo{mbpp/gof.py}
(\code{time\_rescaling\_residuals}, \code{ks\_test\_exp1},
\code{qq\_exp1}), with bucket-level Pearson/dispersion residuals as the
(C4)-regime substitute. \Cref{prop:ch03:nexttype} is the modeling target
of \modrepo{sector\_ranker.py} (\cref{ch:ranking}).
\Cref{prop:ch03:censoring} and the (C2) risk-set rule become validation
checks in the planned \modrepo{data/audits.py} of \cref{ch:data}:
open-spell retention and entry-date consistency are preconditions for the
likelihood to be the one this chapter proved things about.
\end{codehooks}
